\section*{正規部分群と剰余群}

$H$が群$G$の部分群なら, $G/H, H\backslash G$は剰余類の集合として定義された.
$H$が以下で定義する正規部分群であるときには, $G/H$が$H\backslash G$と同一視され,
群の構造が自然に入る.
\begin{dfn*}
	（正規部分群）$H$を群$G$の部分群とする. すべての$g \in G,h \in H$に対し
	$ghg^{-1} \in H$となるとき, $H$を$G$の\textbf{正規部分群}といい,
	$H\lhd G$あるいは$G\rhd H$と書く.
\end{dfn*}

$G$がアーベル群で$H$が任意の部分群なら, $ghg^{-1}=gg^{-1}h=1_Gh=h \in H$
となるので, $H$は正規部分群である.

\begin{tcolorbox}
	\begin{prp*}
		$G_1,G_2$が群で$\phi:G_1 \to G_2$が準同型なら,
		$\mathrm{Ker}(\phi)$は$G_1$の正規部分群である.
	\end{prp*}
\end{tcolorbox}
\begin{proof}
	$g \in G_1, h\in \mathrm{Ker}(\phi)$なら,
	\[ \phi(ghg^{-1})=\phi(g)\phi(h)\phi(g^{-1})
    =\phi(g)1_{G_2}(\phi(g))^{-1}=\phi(g)(\phi(g))^{-1}=1_{G_2}\]
	となるので, $ghg^{-1}\in \mathrm{Ker}(\phi)$となる. 
	よって, $\mathrm{Ker}(\phi)\lhd G_1$である.
\end{proof}\leavevmode

$\mathrm{GL}_n(\mathbb{R}) \in g \mapsto \mathrm{det} g \in \mathbb{R}^\times$
は全射準同型で,$\mathrm{Ker}(\mathrm{det})=\mathrm{SL}_n(\mathbb{R})$であった.
したがって, $\mathrm{SL}_n(\mathbb{R}) \lhd \mathrm{GL}_n(\mathbb{R})$である.
同様に, $\mathrm{SL}_n(\mathbb{C}) \lhd \mathrm{GL}_n(\mathbb{C})$である.
\newline

$G=\mathrm{GL}_2(\mathbb{R})$とし,
\[a(t_1, t_2)=\begin{pmatrix}	t_1 & 0 \\ 0 & t_2 \end{pmatrix},\quad
n(u)=\begin{pmatrix}	1 & 0 \\ u & 1 \end{pmatrix}\]
とおいて
$B=\{a(t_1,t_2)n(u) \mid t_1, t_2 \in \mathbb{R}^\times, u \in \mathbb{R}\},
N=\{n(u)\mid u \in \mathbb{R}\}$とする.
\begin{align*}
  a(t_1,t_2)n(u)n(u')(a(t_1,t_2)n(u))^{-1}
  &=a(t_1,t_2)n(u)n(u')n(-u)a(1/t_1,1/t_2)\\
  &=a(t_1,t_2)n(u')a(1/t_1,1/t_2)\quad (\because n(u)n(u')=n(u')n(u))\\
  &=n(t_2u/t_1)\in N
\end{align*}
となる. したがって$N\lhd B$である.

$B$は$\mathrm{GL}_2(\mathbb{R})$の
ボレル(Borel)部分群と呼ばれるものの一つである。\newline

\begin{tcolorbox}
  \begin{prp*}
    $N$は$G$の部分群で, $G/N$はそれぞれ部分集合$S,T$で生成されているとする.
    このとき, すべての$x\in S, y\in T$に対し$xyx^{-1}, x^{-1}yx \in N$なら,
    $N$は正規部分群である. もし$G$が有限群なら, 
    条件$xyx^{-1}\in N$だけで十分である.
  \end{prp*}
\end{tcolorbox}
\begin{proof}
  まず$S$の任意の元$s$に対し, $xNx^{-1},x^{-1}Nx\subset N$であることを示す.
  $N$の任意の元は$T$の元による語なので, $y_1,y_2,\cdots,y_n \in T$により
  $y_1^{\pm 1}\cdots y_n^{\pm 1}$という形をしている. すると,
  \[ xy_1^{\pm 1}\cdots y_n^{\pm 1}x^{-1}
  =xy_1^{\pm 1}x^{-1} xy_2^{\pm 1}x^{-1} \cdots xy_n^{\pm 1}x^{-1}\]
  だが, $xy_ix^{-1} \in N$なので, 
  $(xy_ix^{-1})^{-1} = xy_i^{-1}x^{-1} \in N$でもある.
  したがって, 上の元は$N$の元となる.
  $x^{-1}Nx \subset N$となることも同様である.

  $G$の任意の元は$S$の元による語である. つまり, $x_1,\cdots x_m \in S$
  があり, $x_1^{\pm 1}\cdots x_m^{\pm 1}$という形をしている.すると,
  \begin{align*}
    gNg^{-1} &= (x_1^{\pm 1}\cdots x_m^{\pm 1})N
    (x_1^{\pm 1}\cdots x_m^{\pm 1})^{-1} \\
     &= (x_1^{\pm 1}\cdots x_{m-1}^{\pm 1})x_m^{\pm 1}N
     (x_m^{\pm 1})^{-1}(x_1^{\pm 1}\cdots x_{m-1}^{\pm 1})^{-1}\\
     & \in (x_1^{\pm 1}\cdots x_{m-1}^{\pm 1})N
     (x_1^{\pm 1}\cdots x_{m-1}^{\pm 1})^{-1}.
  \end{align*}
  $m$に関する帰納法でこれは$N$の元となる. よって, $N$は正規部分群である.
  
  もし$G$が有限群なら, $n=|G|$とすると, 任意の$x\in S$に対し
  $x^n=1_G$である. $x^{-1}=x^{n-1}$なので, $S$の元の語として,
  $x_1,\cdots x_m \in S$により$x_1\cdots x_m$
  という形をしたものだけ考えればよい.したがって, 上の証明より,
  $xyx^{-1} \in N$がすべての$x\in S, y \in T$に対し成り立てば十分である.
\end{proof}
要するに, 正規部分群であることの判定には, 
$G,N$ともに生成元だけ考えればよい.\newline

\begin{tcolorbox}
  $G$を群, $S\subset G$とする. このとき,
  \[ N=\langle \{xyx^{-1} \mid x \in G, y \in S\} \rangle \]
  は$S$を含む最小の$G$の正規部分群である.
\end{tcolorbox}\leavevmode

$G=\mathfrak{G}_3, N=\langle (123) \rangle$とする.
$G$は$S=\{(123),(12)\}$で生成されている.

\quad $(123)(123)(123)^{-1}=(123) \in N,$ 

\quad $(12)(123)(12)^{-1}=(12)(123)(12)=(132)=(123)^2 \in N$

である. $G$は有限群なので, $N$は正規部分群である.\newline

$N$が群$G$の正規部分群で$g \in G$なら, $gN=Ng$である.
\begin{proof}
  $n\in N$なら, $n'=gng^{-1}$とおくと$n'\in N$である.
  よって, $gn=n'g \in Ng$である.
  これがすべての$n \in N$に対して成り立つので, $gN\subset Ng$である.
  同様に$Ng\in gN$も成り立つので, $gN=Ng$である.
\end{proof}
$N$が正規部分群なら, 左剰余類と右剰余類が一致する.

$\pi$を自然な写像$G\to G/N$とする. つまり, $g \in G$に対し
$\pi(g)=gN \in G/N$である.
$G/N$の二つの元を剰余類の代表元$g,h \in G$により$gN,hN$と表す.
このとき, $gN,hN$の積を
\[(gN)(hN)\overset{def}{=}ghN\]
と定義する. この定義が代表元$g,h$の取り方によらないことを示す.
$gN,hN$の任意の元はそれぞれ$n,n'\in N$により$gn,hn$と書ける.
すると, $gnhn'=ghh^{-1}nhn'$だが, $h^{-1}nh \in N$なので,
$h^{-1}nhn' \in N$. よって, $gnhn',gh$の剰余類は等しい.
したがって, 上の定義はwell-definedな写像
\[G/N\times G/N \in (gN,hN) \mapsto ghN \in G/N\]
を定義する.

\begin{tcolorbox}
	\begin{thm*}
		$G/N$は上の演算により群になる.
	\end{thm*}	
\end{tcolorbox}
\begin{proof}
	$1_GN=N$が単位元となることは明らかである.
	$g,h,k \in G$なら, $(gh)k=g(hk)$なので,
	\begin{align*}
		((gN)(hN))(kN) &= (ghN)(kN)=((gh)k)N=(g(hk))N\\
		&=(gN)((hk)N)=(gN)((hN)(kN))
	\end{align*}
	となり, 結合法則が成り立つ. 逆元の存在も同様である.
\end{proof}\leavevmode

\begin{dfn*}
	（剰余群）$G/N$に上の演算を考えたものを, $G$の$N$による商群または剰余群という.
\end{dfn*}

\begin{tcolorbox}
	\begin{prp*}
		自然な写像$\pi: G \to G/N$は群の全射準同型である.
		また$\mathrm{Ker}(\pi)=N$である.
	\end{prp*}
\end{tcolorbox}
\begin{proof}
	$\pi$が全射であることは$G/N$の定義より明らかである.
	$G/N$の積の定義より$\pi$は準同型である. $G/N$の単位元は$N$なので,
	$g\in G$で$\pi(g)=gN=N$であることと$g\in N$であることは同値である.
	よって, $\mathrm{Ker}(\pi)=N$である.
\end{proof}\leavevmode

$n\neq 0$を整数とし, 剰余群$\mathbb{Z}/n\mathbb{Z}$を考える.
$x \in \mathbb{Z}$の$n\mathbb{Z}$に関する剰余類を$\overline{x}$とする.
剰余群$\mathbb{Z}/n\mathbb{Z}$の演算の定義は$x,y \in \mathbb{Z}$に対して
$(x+n\mathbb{Z})+(y+n\mathbb{Z})=(x+y+n\mathbb{Z})$とするものである.
もし$x+y=qn+r$で$q,r\in \mathbb{Z}, 0\leq r < |n|$なら, 
$x+y+n\mathbb{Z}=r+n\mathbb{Z}$である.
$n>0$で$0\leq x,y <n$なら, この定義は$\mathbb{Z}/n\mathbb{Z}$
の和の定義と一致する.

ここでは, $\mathbb{Z}, \mathbb{Z}/n\mathbb{Z}$の
加法による群構造を考えているが, $n\mathbb{Z}\subset \mathbb{Z}$は
$\mathbb{Z}$の「イデアル」と呼ばれるものである.\newline

$G=\mathfrak{S}_3, N=\langle (123) \rangle$とおく.
\[ 6 = |\mathfrak{S}_3| = (\mathfrak{S}_3:N)|N|
=3(\mathfrak{S}_3:N)=3|\mathfrak{S}_3/N|\]
なので, $|\mathfrak{S}_3/N|=2$である.
$(12)\notin N$なので, $(12)N \in \mathfrak{S}_3/N$は単位元ではない.
2は素数なので, $\mathfrak{S}_3/N$は$(12)N$で生成される巡回群である.

\subsubsection*{演習問題}
\textbf{2.8.1}\quad 次の群$G$の部分群$H$が正規部分群であるかどうか判定せよ。
\begin{enumerate}[label=(\arabic*)\;]	
	\item $H=\mathfrak{S}_3 \subset G=\mathfrak{S}_4$.
	\item $H=\mathrm{SO}(2) \subset G=\mathrm{GL}_2(\mathbb{R})$.
	\item $H=\mathrm{GL}_2(\mathbb{R}) \subset G=\mathrm{GL}_2(\mathbb{C})$.
	\item $H=\{1,(12)(34),(13)(24),(14)(23)\} \in G=\mathfrak{S}_4$.
	\item $G$は$\mathrm{GL}_2(\mathbb{R})$の元で下三角行列であるもの全体よりなる群、
			$H$は$G$の元で対角成分が等しい元よりなる部分群.
\end{enumerate}
\par
\begin{enumerate}[label=(\arabic*)\;]
	\item $H=\{1,(123),(132),(12),(23),(13)\}$であるが,
		$(14) \in G$に対して$(14)(12)(14)^{-1}=(124)(14)=(24) \notin H$
		となる. したがって正規部分群ではない.
	\item $h=\begin{pmatrix}0&1\\1&0\end{pmatrix}\in H,
		g=\begin{pmatrix}1&0\\1&1\end{pmatrix}\in G$に対して,
		\[ghg^{-1}=\begin{pmatrix}1&0\\1&0\end{pmatrix}
		\begin{pmatrix}0&1\\1&0\end{pmatrix}
		\begin{pmatrix}1&0\\-1&1\end{pmatrix}
		=\begin{pmatrix}0&1\\1&1\end{pmatrix}
		\begin{pmatrix}1&0\\-1&1\end{pmatrix}
		=\begin{pmatrix}-1&1\\0&1\end{pmatrix} \notin H\]
		なので,正規部分群ではない.

		($\mathrm{O}(2)=\{A\in \mathrm{GL}_2(\mathbb{R})\mid {}^TAA=I_2)\},
		\mathrm{SO}(2)=\{A \in \mathrm{O}(2)\mid \det A=1\}$)
	\item $\sqrt{-1}=i$とする.
		$h=\begin{pmatrix}1&0\\1&1\end{pmatrix}\in H,
		g=\begin{pmatrix}1&0\\0&i\end{pmatrix}\in G$に対して,
		\[ghg^{-1}=\begin{pmatrix}1&0\\0&i\end{pmatrix}
		\begin{pmatrix}1&0\\1&1\end{pmatrix}
		\begin{pmatrix}1&0\\0&-i\end{pmatrix}
		=\begin{pmatrix}1&0\\i&i\end{pmatrix}
		\begin{pmatrix}1&0\\0&-i\end{pmatrix}
		=\begin{pmatrix}1&0\\i&1\end{pmatrix} \notin H\]
		なので,正規部分群ではない.

\end{enumerate}

\textbf{2.8.2}\quad $H$を群$G$の指数2の部分群とする. 
このとき, $H$は$G$の正規部分群であることを証明せよ.

\textbf{2.8.3}\quad $N_1, N_2$が群$G$の正規部分群なら, 
$N_1 N_2$も$G$の正規部分群であることを証明せよ.

\begin{proof}
	$n_1 \in N_1, n_2 \in N_2$とする.
	$N_1, N_2$は$G$の部分群なので$1_G\in N_1$かつ$1_G \in N_2$.
	任意の$n_1, n_2$に対して$n_1n_2 \in N_1N_2$であり,
	$n_2=1_G$とすると$n_1n_2=n_11_G=n_1 \in N_1N_2$.
	同様にして$n_2 \in N_1N_2$である.
	したがって$n_1, n_2 \in N_1N_2$であり, 
	これらは積について閉じているので$N_1N_2$は$G$の部分群となる.

	$N_1, N_2$は$G$の正規部分群なので,$g \in G$に対して
	$gn_1g^{-1} \in N_1, gn_2g^{-1} \in N_2$である.
	よって$(gn_1g^{-1})(gn_2g^{-1})=gn_1(g^{-1}g)n_2g^{-1}
	=g(n_1n_2)g^{-1} \in N_1N_2$となるので,
	$N_1N_2$は$G$の正規部分群である.
\end{proof}

\textbf{2.8.4}\quad $\mathfrak{S}_3$の部分群をすべて求めよ.
そのなかで正規部分群はどれか?

\textbf{2.8.5}\quad 四元数群の部分群をすべて求めよ.
そのなかで正規部分群はどれか?
